% PDF-to-TeX transcription by the assistant; not the original downloaded file.

\begin{theorem}[38.11, Borsuk--Ulam]
Think of $S^n$ as the unit vectors in $\R^{n+1}$. For any continuous function $f:S^n\to\R^n$, there exists $x\in S^n$ such that $f(x)=f(-x)$.
\end{theorem}
\begin{proof}
Suppose that no such $x$ exists. Then we may define a continuous function $g:S^n\to S^{n-1}$ by
\[
g:x\longmapsto\frac{f(x)-f(-x)}{\|f(x)-f(-x)\|}.
\]
Note that $g(-x)=-g(x)$: $g$ is equivariant with respect to the antipodal action. It descends to a map $\bar g:\mathbf{RP}^n\to\mathbf{RP}^{n-1}$.

We claim that $\bar g_*:H_1(\mathbf{RP}^n)\to H_1(\mathbf{RP}^{n-1})$ is nontrivial. To see this, pick a basepoint $b\in S^n$ and choose a $1$-simplex $\sigma:\Delta^1\to S^n$ such that $\sigma(e_0)=b$ and $\sigma(e_1)=-b$. The group $H_1(\mathbf{RP}^n)$ is generated by the cycle $p\sigma$. The image of this cycle in $H_1(\mathbf{RP}^{n-1})$ is represented by the loop $\bar g p\sigma$ at $\bar b=pb$, which is the image of the $1$-simplex $g\sigma$ joining $gb$ to $g(-b)=-g(b)$. The class of this $1$-simplex thus generates $H_1(\mathbf{RP}^{n-1})$.

Therefore $\bar g$ is nontrivial in $H_1(-;\mathbb F_2)$, and hence also in $H^1(-;\mathbb F_2)$. Writing $a_n$ for the generator of $H^1(\mathbf{RP}^n;\mathbb F_2)$, we must have $a_n=\bar g^*a_{n-1}$, and consequently
\[
a_n^n=(\bar g^*a_{n-1})^n=\bar g^*(a_{n-1}^n).
\]
But $H^n(\mathbf{RP}^{n-1};\mathbb F_2)=0$, so $a_{n-1}^n=0$; while $a_n^n\neq0$. This is a contradiction.
\end{proof}

